
% ════════════════════════════════════════
% 模型假设（03.模型假设.tex）
% ════════════════════════════════════════
\section{模型假设}

为简化问题，本文做出以下假设：

\begin{itemize}[itemindent=2em]
\item \textbf{假设1：}各 10 分钟间隔内的功率保持恒定，购电量、充放电量与功率成正比，忽略更短时间尺度上的功率波动。
\item \textbf{假设2：}储能设备的充放电效率均为 90\%，容量上下限与充放电功率限值在整个运行周期内保持不变，忽略电池老化与自放电。
\item \textbf{假设3：}每天 0:00 制定计划时，当天各间隔的小区负载与光伏出力均尚未发生，不能直接引用其实际值，只能依据历史数据与滚动预报建立预测模型给出估计，预测必然存在误差；实际偏差通过紧急购电机制兜底。
\item \textbf{假设4：}外网购电功率无上限，微网与外部电网可随时完成任意规模的购电交易，交易电价仅取决于对应时段的电价数据。
\item \textbf{假设5：}各日 0:00 依据当日预测独立决策，储电量跨日传递；日末储电量不作终端约束，每日末端可放至下限，次日谷时再充，这是逐日滚动下的经济性选择。
\item \textbf{假设6：}实际运行时储能严格按 0:00 制定的充放电计划执行，真实负载与光伏到达后，出现的供电缺口全部由 5 倍电价的紧急购电补足，富余电量计为弃电。
\end{itemize}
